\chapter{协作机器人碰撞安全：检测、ISO/TS 15066 与碰撞后策略}\label{ch:arm-collision-safety}

% ════════════════════════════════════════════════════════════════
%  本章：把机械臂安全工程从「能写出阻抗/力控律」推进到「能让一台与人共处的机械臂
%  在碰撞时不伤人、且事后能安全收场」。检测理论(动量观测器)与无源性已在操作空间章
%  讲透，本章只做安全工程落地：阈值/方向分离、ISO/TS 15066 力压阈值、有效质量限速、
%  碰撞后 Reflex/Retract/Comply、能量分级响应。源：Tutorial F06(碰撞安全部分)，
%  转写为本书记号(右扰动、tau=Jᵀf、Λ=(JM⁻¹Jᵀ)⁻¹)。
% ════════════════════════════════════════════════════════════════

\begin{practice}[前置自测]
开始之前，先试答下列问题。若已能流畅作答，可只速读本章的阈值表与速查卡；若卡住，正是本章要为你解决的。
\begin{enumerate}
  \item 一台 $18$\,kg 的机械臂以 $0.5$\,m/s 撞到人手，碰撞峰值力由“机器人的总质量”决定吗？如果不是，由什么决定？
  \item “我把阻抗刚度 $\mathbf K_d$ 调得很低、误差限幅 $e_{\max}$ 也很小，准静态接触力只有 $25$\,N，远低于人体阈值”——这是否就自动满足了协作安全标准？哪几种情形会让它失效？
  \item 碰撞被检测到之后，让机械臂\emph{立即刹停并保持位置}（Reflex），听起来最安全。为什么在“人被夹在机器人与桌面之间”的场景里，它反而是三种策略里最危险的？
  \item 在销孔装配的插入阶段，沿插入方向 $z$ 的接触力可以高达 $30$\,N（正常工艺力），而侧向 $1$\,N 的异常力却意味着卡死。用一个固定的力阈值能同时“放过工艺力”又“抓住卡死”吗？
  \item ISO/TS 15066 给了“准静态”与“瞬态”两套力阈值，后者约为前者的两倍。这个“两倍”是工程拍脑袋的安全系数，还是有生物力学依据？
\end{enumerate}
\end{practice}

\section*{本章目标}
学完本章，你应当能够：
\begin{enumerate}
  \item \textbf{区分} ISO/TS 15066 的两类协作模式——速度与间隔监控（SSM）与功率/力限制（PFL）——并说清力控只在 PFL 模式下承担安全红线；
  \item \textbf{设置}基于动量观测器残差 $\mathbf r$ 的碰撞检测阈值，并用\emph{方向分离} $r_{\rm op}/r_\perp$ 在“放过工艺力”与“抓住异常力”之间取得灵敏度；
  \item \textbf{读懂并使用} ISO/TS 15066 PFL 的力/压力阈值体系（按身体区域、区分准静态与瞬态），并解释“瞬态$\approx 2\times$准静态”的生物力学来由；
  \item \textbf{计算}构型相关的有效碰撞质量 $m_{\rm eff}=1/(\mathbf n^\top\boldsymbol\Lambda^{-1}\mathbf n)$ 与由此导出的碰撞速度上界 $v_{\max}$，并说明它\emph{为何不等于}自动满足 PFL；
  \item \textbf{设计}碰撞后恢复策略（Reflex / Retract / Comply），论证“准静态合规”为何把 Comply/Retract 推为首选、而 Reflex 的残余夹持力是一个安全悖论；
  \item \textbf{构建}以能量为度量的分级响应（Level\,0--3），把检测、限速与碰撞后策略接成一条运行时安全链。
\end{enumerate}

\subsection*{本章知识导航}
本章是一条“\emph{从感知到合规、再到事后收场}”的安全工程主线。结构如下：

\begin{center}
\small
\begin{tikzpicture}[
  node distance=4mm and 7mm, >=Stealth,
  blk/.style={draw, rounded corners, align=center, font=\small, inner sep=3pt, fill=main!6, draw=main!60!black},
  grp/.style={draw, rounded corners, align=center, font=\small, inner sep=3pt, fill=second!6, draw=second!60!black},
  saf/.style={draw, rounded corners, align=center, font=\small, inner sep=3pt, fill=third!8, draw=third!60!black},
  ln/.style={->, thick, gray!60!black}]
\node[grp] (mode) {协作模式\\ SSM / PFL};
\node[blk, right=of mode] (det) {碰撞检测\\ 残差 $\mathbf r$ + 阈值};
\node[blk, right=of det] (dir) {方向分离\\ $r_{\rm op}/r_\perp$};
\node[saf, below=8mm of mode] (iso) {ISO/TS 15066\\ 力/压力阈值};
\node[saf, right=of iso] (mass) {有效质量\\ $m_{\rm eff}$, $v_{\max}$};
\node[saf, right=of mass] (post) {碰撞后策略\\ Reflex/Retract/Comply};
\node[blk, below=8mm of iso, fill=third!8, draw=third!60!black] (level) {能量分级响应\\ Level 0--3};
\draw[ln] (mode)--(det); \draw[ln] (det)--(dir);
\draw[ln] (mode.south)--(iso.north);
\draw[ln] (iso)--(mass); \draw[ln] (mass)--(post);
\draw[ln] (det.south) -- ++(0,-4mm) -| (mass.north);
\draw[ln] (iso.south) -- (level.north);
\draw[ln] (post.south) |- (level.east);
\end{tikzpicture}
\end{center}

\noindent\textbf{推荐路径}：第~\ref{sec:cs-motivation}~节确立“法规即红线”的动机与两种模式的分工，是全章的坐标系；第~\ref{sec:cs-detection}~节把检测落到\emph{阈值与方向分离}（检测\emph{原理}已在操作空间章给出，本节只做工程化）；第~\ref{sec:cs-iso}~节是安全合规的核心数据与判据；第~\ref{sec:cs-mass}~节把“速度多快才安全”量化为构型相关的不等式，并刻画它与阻抗参数的关系及限定；第~\ref{sec:cs-postcollision}~节解决“撞到之后做什么”，第~\ref{sec:cs-energy}~节把前面接成运行时安全链。时间有限可走 \ref{sec:cs-motivation}$\to$\ref{sec:cs-iso}$\to$\ref{sec:cs-postcollision} 的“合规—收场”最短路径。

\subsection*{前置知识桥接}
本章把三块前置能力焊到“安全工程”上。其一，碰撞\emph{检测}的理论核心——De Luca--Mattone 广义动量观测器——已在 \cref{thm:os-observer} 完整推导（残差 $\mathbf r$ 以一阶滤波收敛到外力矩 $\boldsymbol\tau_{\rm ext}$，无需测关节加速度），本章直接以 $\mathbf r$ 为输入，不重复推导，只解决\emph{阈值怎么定、方向怎么分}；\cref{ch:force-wbc} 亦从全身控制视角给出残差式碰撞检测。其二，碰撞动力学的“感觉重量”——构型相关的有效质量 $m_{\rm eff}=1/(\mathbf n^\top\boldsymbol\Lambda^{-1}\mathbf n)$（\cref{eq:os-meff}）与操作空间惯量 $\boldsymbol\Lambda=(\mathbf J\mathbf M^{-1}\mathbf J^\top)^{-1}$（\cref{ch:arm-dyn}）——是本章一切速度上界的基石。其三，能量端口 $P=\mathbf f^\top\mathbf v$（\cref{eq:os-port-power}）与静力对偶 $\boldsymbol\tau=\mathbf J^\top\mathbf f$ 提供了把“安全”记成“能量账”的语言。本章是 \cref{ch:arm-prac} 之前的安全前置：任何要让机械臂与人共处的实践，都得先过本章这一关。

\subsection*{如果跳过本章会怎样}
\begin{itemize}
  \item 你会把“通过了 SLAM/控制的精度指标”误当成“可以部署到人身边”，而不知道协作部署的真正闸门是 ISO/TS 15066 的力/压力限值——产品无法合规上市；
  \item 你会给碰撞检测设一个全局固定阈值，于是要么在正常工艺力下频繁误触发、任务做不下去，要么把阈值抬高到漏掉真实碰撞；
  \item 你会本能地选“撞到就急停并锁住”，却在“人被夹住”的最危险场景里让残余力持续施加在人身上——一个看似最安全、实则违规的设计。
\end{itemize}

\begin{note}
\textbf{本章记号与约定.}\;
关节空间动力学沿用全书 $\mathbf M(\mathbf q)\ddot{\mathbf q}+\mathbf C(\mathbf q,\dot{\mathbf q})\dot{\mathbf q}+\mathbf g(\mathbf q)=\boldsymbol\tau+\boldsymbol\tau_{\rm ext}$，外部接触力矩 $\boldsymbol\tau_{\rm ext}$ 显式在右端；末端 wrench $\mathbf f$ 与关节力矩经静力对偶 $\boldsymbol\tau=\mathbf J^\top\mathbf f$ 相连。动量观测器残差记 $\mathbf r$（\cref{thm:os-observer}，增益 $\mathbf K_I$、时间常数 $\mathbf K_I^{-1}$），它估计 $\boldsymbol\tau_{\rm ext}$。操作空间惯量 $\boldsymbol\Lambda=(\mathbf J\mathbf M^{-1}\mathbf J^\top)^{-1}$，有效碰撞质量沿 \cref{eq:os-meff}。能量端口功率 $P=\mathbf f^\top\mathbf v$（\cref{eq:os-port-power}）。本章\emph{不}重新约定阻抗记号，沿用控制卷 $\mathbf K_d,\mathbf D_d$；接触力压阈值用 $F,p$，下标 $\rm qs$（quasi-static，准静态）与 $\rm tr$（transient，瞬态）。\textbf{数值阈值多来自 ISO/TS 15066:2016 的 Mainz 生物力学数据}，本章给出代表性数值并对不确定处打 \pz 注明，权威值以现行标准正文为准。
\end{note}

% ════════════════════════════════════════════════════════════════
\section{动机：当法规成为力控的红线}
\label{sec:cs-motivation}

\paragraph{动机：精度达标 $\neq$ 可以部署到人身边。}
前面若干章把机械臂的运动学、动力学、轨迹与力/阻抗控制讲透，衡量标准始终是\emph{任务品质}——更小的跟踪误差、更稳的接触力。但当这台机械臂被放到\emph{人身边}（协作机器人，cobot），衡量标准多了一条、且优先级最高：\textbf{它撞到人时不能造成伤害}。这不再是性能问题，而是\emph{合规}问题——由安全标准给出一条任何控制器都不得越过的红线。理解这条红线，是把“会写力控”升级为“能交付协作产品”的分水岭。

\paragraph{两种协作模式：谁来守红线。}
ISO 10218 与其技术规范 ISO/TS 15066 把人机协作落地为若干\emph{协作模式}，与本书最相关的是两种互补范式（\cref{fig:cs-modes}）：
\begin{itemize}
  \item \textbf{SSM（Speed and Separation Monitoring，速度与间隔监控）}：用外部感知（安全激光扫描、视觉）实时测量人与机器人的间距，间距越近、允许速度越低，触及保护性间隔则停机。安全由\emph{感知 + 速度规划}保障，碰撞被\emph{预防}在发生之前——理想情况下根本不接触。
  \item \textbf{PFL（Power and Force Limiting，功率与力限制）}：允许人与机器人在同一空间\emph{同时}作业、甚至允许接触，但要求一旦碰撞，接触的\emph{力}与\emph{压力}不超过人体可承受阈值。安全由\emph{机器人本体}保障——轻量化机械结构、力矩限制、以及\textbf{力控器对接触力的主动管理}。
\end{itemize}
关键分工：\textbf{SSM 模式下安全主要靠感知与限速，力控不直接守红线；PFL 模式下力控是红线的直接执行者}。本章聚焦 PFL——因为它把“安全”变成了一个关于力、压力、有效质量与碰撞后策略的\emph{控制工程}问题。

\begin{figure}[ht]
\centering
\begin{tikzpicture}[>=Stealth, font=\small, scale=1.0]
  % SSM panel
  \node[font=\bfseries] at (-3.1,2.0) {SSM：预防接触};
  \draw[main!60!black, thick] (-4.6,0) rectangle (-4.2,1.4);  % robot base
  \node[main!60!black] at (-4.4,-0.3) {机器人};
  \fill[second!70!black] (-1.7,0.7) circle (0.18);            % human
  \node[second!70!black] at (-1.7,-0.3) {人};
  \draw[<->, dashed, third!70!black, thick] (-4.1,0.9) -- (-1.95,0.7) node[midway, above]{间距 $d$};
  \draw[->, gray] (-4.1,1.6) .. controls (-3.2,2.0) .. (-2.0,1.4) node[right,gray]{$v\!\downarrow$ 随 $d\!\downarrow$};
  % divider
  \draw[gray!50] (0,-0.6) -- (0,2.2);
  % PFL panel
  \node[font=\bfseries] at (3.1,2.0) {PFL：限制接触力};
  \draw[main!60!black, thick] (1.4,0) rectangle (1.8,1.4);
  \node[main!60!black] at (1.6,-0.3) {机器人};
  \fill[second!70!black] (3.4,0.85) circle (0.18);
  \node[second!70!black] at (3.4,-0.3) {人};
  \draw[<->, red!60!black, very thick] (1.85,0.9) -- (3.2,0.86) node[midway, above]{$F\le F_{\max}$};
  % compliant-contact spring drawn as plain zigzag (no decorations.pathmorphing dependency)
  \draw[red!55!black] (1.9,0.55) -- (2.05,0.65) -- (2.25,0.45) -- (2.45,0.65) -- (2.65,0.45)
        -- (2.85,0.65) -- (3.05,0.45) -- (3.18,0.55);
\end{tikzpicture}
\caption{两种协作模式的分工。左：SSM 用间距 $d$ 调制速度 $v$，把碰撞\emph{预防}在发生前，安全靠感知与限速；右：PFL 允许接触，但用本体（轻量结构 + 力控）把接触力压在 $F_{\max}$ 内，安全靠力控守红线。本章聚焦 PFL。}
\label{fig:cs-modes}
\end{figure}

\begin{insight}
为什么力控章要专门辟出“安全”一节？因为 PFL 把一个\emph{标准条文}翻译成了一组\emph{控制约束}：碰撞力上界 $\to$ 速度上界 $\to$ 阻抗参数上界 $\to$ 碰撞后必须主动卸力。这条翻译链上的每一环，都是前面章节学过的力控工具的直接应用——安全不是另起炉灶，而是力控品质在“以人为约束”下的极限形态。
\end{insight}

% ════════════════════════════════════════════════════════════════
\section{碰撞检测的工程化：阈值与方向分离}
\label{sec:cs-detection}

\paragraph{从“能估外力”到“能判碰撞”。}
\cref{thm:os-observer} 已给出无传感外力估计：广义动量观测器输出残差 $\mathbf r$，它以一阶线性滤波 $\dot{\mathbf r}=\mathbf K_I(\boldsymbol\tau_{\rm ext}-\mathbf r)$ 指数收敛到真实外部力矩 $\boldsymbol\tau_{\rm ext}$，且只用可直接测量的 $\mathbf q,\dot{\mathbf q},\boldsymbol\tau$。检测\emph{原理}至此完备，本节只回答两个工程问题：\textbf{阈值定多大}、以及当“正常接触”与“碰撞”混在一起时\textbf{怎么把它们分开}。

\paragraph{阈值选择：灵敏度与误报的拉锯。}
最朴素的检测是对残差范数设一个固定阈值 $\|\mathbf r\|>r_{\rm th}\Rightarrow$ 触发安全响应。困难在于 $r_{\rm th}$ 的两难（\cref{fig:cs-threshold}）：
\begin{itemize}
  \item $r_{\rm th}$ \emph{太低}：模型误差（关节摩擦、传动柔性、负载标定偏差）与正常操作力都会把 $\|\mathbf r\|$ 顶过阈值，造成\emph{误触发}，任务被频繁打断；
  \item $r_{\rm th}$ \emph{太高}：真实碰撞的残差要爬升更久才越线，检测\emph{延迟}增大甚至\emph{漏报}，而碰撞峰值力随接触持续时间与穿透深度增长——延迟直接恶化伤害。
\end{itemize}
阈值的下限由\emph{残差噪声底}决定：$\mathbf r$ 的稳态波动来自动力学模型残差与传感噪声，工程上取 $r_{\rm th}$ 略高于无碰撞时 $\|\mathbf r\|$ 的统计上界（如均值加若干倍标准差）。观测器带宽（由增益 $\mathbf K_I$ 设定）则决定了响应速度与噪声放大的折中：增益越大，$\mathbf r$ 越快逼近 $\boldsymbol\tau_{\rm ext}$、检测越快，但噪声底也越高、被迫抬高 $r_{\rm th}$。\pz{存疑：常见工程取值——碰撞检测阈值约 $20$--$30$\,N（折算到末端 wrench）、观测器增益对应带宽数 Hz 到十余 Hz——为 Tutorial 给出的经验范围，随机器人与标定质量变化，部署前应实测残差噪声底重新标定，不应直接套用。}

\begin{figure}[ht]
\centering
\begin{tikzpicture}[>=Stealth, font=\small, scale=1.0]
  \draw[->] (0,0) -- (7.4,0) node[right]{$t$};
  \draw[->] (0,0) -- (0,3.1) node[above]{$\|\mathbf r\|$};
  % noisy baseline
  \draw[gray!70] (0.1,0.45) .. controls (0.9,0.6) and (1.4,0.32) .. (2.0,0.5)
        .. controls (2.6,0.66) and (3.2,0.38) .. (3.8,0.52)
        .. controls (4.2,0.6) and (4.5,0.45) .. (4.7,0.52);
  \node[gray] at (1.4,0.95) {残差噪声底};
  % collision rise
  \draw[main!70!black, very thick] (4.7,0.52) .. controls (5.0,0.9) and (5.2,2.1) .. (5.6,2.7);
  \node[main!70!black] at (6.3,2.6) {碰撞发生};
  % low threshold (false trigger)
  \draw[red!60!black, dashed] (0,0.75) -- (7.0,0.75) node[right]{$r_{\rm th}$ 太低};
  \fill[red!60!black] (2.0,0.5) circle (0pt);
  \node[red!60!black, font=\scriptsize] at (2.7,0.78) {误触发$\times$};
  % good threshold
  \draw[third!70!black, thick] (0,1.25) -- (7.0,1.25) node[right]{$r_{\rm th}$ 合适};
  % high threshold (delay)
  \draw[second!70!black, dashed] (0,2.25) -- (7.0,2.25) node[right]{$r_{\rm th}$ 太高};
  \draw[<->, second!70!black, font=\scriptsize] (5.15,1.27) -- (5.42,2.23) node[midway,right]{检测延迟};
\end{tikzpicture}
\caption{碰撞检测阈值的两难。阈值太低（红）被残差噪声底顶穿、误触发；太高（橙）则真实碰撞需爬升更久才越线、检测延迟增大。合适阈值（蓝）略高于噪声底统计上界，兼顾灵敏与抗扰。}
\label{fig:cs-threshold}
\end{figure}

\paragraph{方向分离：让阈值“认得方向”。}
固定的标量阈值有一个根本缺陷：它把\emph{所有方向}的外力一视同仁。但在\emph{接触任务}中，沿操作方向的力往往是\textbf{期望的工艺力}（如销孔装配中沿插入方向 $z$ 的推力可达数十牛），而垂直于操作方向的侧向力才是\emph{碰撞/卡死}的信号。用单一阈值无法两全：抬高它放过工艺力，就漏掉侧向异常；压低它抓住侧向，工艺力就频繁误触发。

解法是把残差\emph{投影}到“操作方向”与“正交补”两个子空间，分设阈值。设末端外力估计 $\hat{\mathbf f}_{\rm ext}$ 由残差经静力对偶反解（\cref{thm:os-observer} 的笛卡尔投影，$\boldsymbol\tau_{\rm ext}=\mathbf J^\top\mathbf f_{\rm ext}$），$\hat{\mathbf n}_{\rm op}$ 为单位操作方向，则
\begin{align}
  r_{\rm op}&=\hat{\mathbf n}_{\rm op}^\top\hat{\mathbf f}_{\rm ext}, &
  r_\perp&=\bigl\|\hat{\mathbf f}_{\rm ext}-r_{\rm op}\,\hat{\mathbf n}_{\rm op}\bigr\|.
  \label{eq:cs-dirsplit}
\end{align}
对操作分量 $r_{\rm op}$ 用\emph{高}阈值（容许工艺力），对正交分量 $r_\perp$ 用\emph{低}阈值（灵敏捕捉侧向碰撞）。这样同一套残差既“放过”了沿工艺方向的大力，又“抓住”了侧向的小异常。\pz{存疑：上式将关节残差 $\mathbf r$ 经 $(\mathbf J^\top)^+$ 反投影为末端 wrench 再做方向分解；当 $\mathbf J$ 接近奇异时反投影病态，工程上须用阻尼最小二乘（操作空间章给出的 $\bar{\mathbf J}$ 动力学一致伪逆是更稳健的选择），方向分离的阈值也宜随构型自适应——细节随平台而定。}

\begin{insight}
方向分离把“灵敏度”从一个标量变成了一个\emph{随任务方向各向异性}的量。它与 \cref{eq:os-meff} 的“有效质量随方向变化”是同一种几何直觉的两面：碰撞的\emph{危险性}（有效质量）与\emph{可检测性}（残差阈值）都依赖方向，安全设计因此天然是“构型相关”的——固定标量参数注定在某些构型/方向上要么过保守、要么不安全。
\end{insight}

\begin{pitfall}[把“检测到碰撞”等同于“知道撞在哪条连杆”]
残差 $\mathbf r$ 给出的是\emph{广义}外力矩，能判断“有外力”、并经雅可比反投影估计末端 wrench；但要进一步\emph{定位}碰撞发生在哪一节连杆（碰撞隔离，collision isolation），需要逐连杆的雅可比与接触模型，是更强的推断。本章的方向分离解决的是“工艺力 vs.\ 侧向异常”的\emph{方向}问题，不等于解决了“撞在哪里”的\emph{位置}问题——别把二者混为一谈。
\end{pitfall}

% ════════════════════════════════════════════════════════════════
\section{ISO/TS 15066：PFL 的力与压力阈值}
\label{sec:cs-iso}

\paragraph{标准的定位：技术规范，而非正式标准。}
ISO/TS 15066:2016 是协作操作安全的\textbf{技术规范}（Technical Specification），\emph{不是}正式国际标准（International Standard）——它是对 ISO 10218-1/2（工业机器人安全）的补充，为协作作业提供具体的安全参数。尽管法律地位是“TS”，它在实践中被协作机器人行业广泛遵循，其力/压力限值已成为产品认证的\emph{事实基准}。对力控设计者，理解 15066 不是可选项：它定义了控制器\textbf{绝对不能超过}的力与压力。

\paragraph{两类碰撞：准静态与瞬态。}
PFL 用\emph{力}与\emph{压力}两个指标评估碰撞严重度，并把碰撞二分为两类——这一二分是整个阈值体系的骨架（\cref{fig:cs-qs-tr}）：
\begin{itemize}
  \item \textbf{准静态接触}（quasi-static，下标 $\rm qs$）：人被\emph{夹}在机器人与一个固定物体（如桌面、墙）之间，接触力\emph{持续}作用、无法卸去。这是最危险的场景，也是 15066 设计的核心目标。
  \item \textbf{瞬态接触}（transient，下标 $\rm tr$）：机器人撞到人后立即弹开或停止，接触力\emph{仅短暂}作用（典型 $<0.5$\,s 之后即卸去）。
\end{itemize}

\begin{figure}[ht]
\centering
\begin{tikzpicture}[>=Stealth, font=\small, scale=1.0]
  % quasi-static: clamping
  \node[font=\bfseries] at (-3.0,2.3) {准静态：夹持};
  \draw[gray!60!black, very thick] (-4.6,-0.1) -- (-1.4,-0.1);   % fixed surface
  \node[gray] at (-4.2,-0.4) {固定面};
  \fill[second!70!black] (-3.0,0.45) circle (0.22);             % human part
  \draw[main!60!black, very thick] (-3.25,1.6) rectangle (-2.75,1.0); % robot
  \draw[->, red!60!black, very thick] (-3.0,0.95) -- (-3.0,0.7);
  \node[red!60!black, font=\scriptsize] at (-2.0,0.55) {$F$ 持续};
  % time profile QS
  \draw[->] (-4.7,-1.6) -- (-1.5,-1.6) node[right,font=\scriptsize]{$t$};
  \draw[->] (-4.7,-1.6) -- (-4.7,-0.7) node[above,font=\scriptsize]{$F$};
  \draw[red!60!black, thick] (-4.6,-1.55) .. controls (-4.3,-0.95) .. (-3.9,-0.9) -- (-1.7,-0.9);
  \node[font=\scriptsize, red!60!black] at (-2.9,-0.7) {不卸去};
  % divider
  \draw[gray!50] (-0.3,-1.7) -- (-0.3,2.5);
  % transient: free impact
  \node[font=\bfseries] at (3.0,2.3) {瞬态：自由碰撞};
  \fill[second!70!black] (3.0,0.45) circle (0.22);
  \draw[main!60!black, very thick] (2.75,1.6) rectangle (3.25,1.0);
  \draw[->, red!60!black, very thick] (3.0,0.95) -- (3.0,0.72);
  \draw[->, third!70!black, thick] (3.3,1.3) .. controls (3.9,1.5) .. (4.2,1.1) node[right,font=\scriptsize]{弹开};
  % time profile TR
  \draw[->] (1.5,-1.6) -- (4.7,-1.6) node[right,font=\scriptsize]{$t$};
  \draw[->] (1.5,-1.6) -- (1.5,-0.7) node[above,font=\scriptsize]{$F$};
  \draw[red!60!black, thick] (1.6,-1.55) .. controls (2.3,-1.55) and (2.6,-0.75) .. (2.9,-0.78)
        .. controls (3.2,-0.8) and (3.5,-1.55) .. (4.4,-1.55);
  \node[font=\scriptsize, red!60!black] at (3.0,-0.6) {$<0.5$\,s};
\end{tikzpicture}
\caption{准静态（夹持，左）与瞬态（自由碰撞，右）。准静态接触力持续不卸、最危险；瞬态力仅短暂作用，故 15066 给瞬态更宽的阈值（约 $2\times$）。}
\label{fig:cs-qs-tr}
\end{figure}

\paragraph{阈值体系：按身体区域分别规定。}
15066 把人体划成约 $29$ 个身体区域，对每个区域分别给出准静态/瞬态的\emph{力}与\emph{压力}阈值（标准正文 Annex A）。\cref{tab:cs-iso} 摘录若干代表性区域的数值，用以建立量级直觉——\textbf{头面部最敏感（阈值低）、四肢与躯干外侧较耐受（阈值高）}。

\begin{table}[H]\centering\small
\caption{ISO/TS 15066:2016 Annex A 表 A.2 PFL 力/压力阈值（代表性区域摘录，建立量级直觉用）}
\label{tab:cs-iso}
\begin{tabular}{lcccc}
\toprule
身体区域（取样点） & 准静态力 $F_{\rm qs}$ [N] & 准静态压力 $p_{\rm qs}$ [N/cm$^2$] & 瞬态力 $F_{\rm tr}$ [N] & 瞬态压力 $p_{\rm tr}$ [N/cm$^2$] \\
\midrule
前额（额中）       & 130 & 130 & ---$^\dagger$ & ---$^\dagger$ \\
颞部              & 130 & 110 & ---$^\dagger$ & ---$^\dagger$ \\
胸骨              & 140 & 120 & 280 & 240 \\
腹部（腹肌）       & 110 & 140 & 220 & 280 \\
骨盆（盆骨）       & 180 & 210 & 360 & 420 \\
上臂（三角肌）     & 150 & 190 & 300 & 380 \\
前臂（前臂肌）     & 160 & 180 & 320 & 360 \\
手背              & 140 & 200 & 280 & 400 \\
食指指腹           & 140 & 300 & 280 & 600 \\
大腿（股肌）       & 220 & 250 & 440 & 500 \\
\bottomrule
\end{tabular}
\end{table}

\noindent{\footnotesize $^\dagger$ 标准把颅/额/面列为\emph{临界区}，接触\emph{不被允许}，故\textbf{不给瞬态值}（“not applicable”）——瞬态$2\times$放大规则\emph{不}适用于该三区。表 A.2 中“力”按 12 个身体区域给出、“压力”按 29 个具体取样点给出，故同一行的力与压力可能取自不同分辨率（本表力列对应所在身体区域、压力列对应括注取样点）。}

\begin{finenote}
\pz{存疑：\cref{tab:cs-iso} 已据 ISO/TS 15066:2016 Annex A 表 A.2 正文逐值校正（准静态力按身体区域、准静态压力按具体取样点；瞬态$=2\times$准静态，颅/额/面除外），仅摘录代表性区域用以建立量级直觉。标准修订版可能调整数值与人群覆盖；安全设计须以购得的现行标准正文为唯一权威来源，切勿以本表为认证依据。}
\end{finenote}

\paragraph{“瞬态$\approx 2\times$准静态”的生物力学依据。}
注意 \cref{tab:cs-iso} 中绝大多数区域的瞬态阈值恰为准静态的两倍——标准明确规定，瞬态接触的力/压力限值由准静态值乘以一个\textbf{瞬态倍率（取 $2$）}得到（表 A.2 的 $P_T,F_T$ 列），\emph{唯有颅/额/面三个临界区例外}：那里接触本就不被允许，故不给瞬态值、不施倍率。这个 $2$ \emph{不是}随手取的安全系数，而有生物力学根据：人体组织对\emph{短时}载荷的疼痛/损伤耐受高于\emph{持续}载荷——短暂冲击在组织产生不可逆损伤前即已卸去；标准称“瞬态限值\emph{至少}可达准静态的两倍”，取 $2$ 是\emph{保守}下界而非精确折算。准静态阈值本身源自德国美因茨（Mainz）大学受 ISO 委员会委托做的疼痛阈值实验：$100$ 名健康成年受试者、$29$ 个具体取样点，压力取所记录值的 $75$ 百分位为“疼痛起始”阈；而“力”阈值另由一项独立研究（汇集 $188$ 处文献来源、以 AIS\,1 轻伤为界）给出，故力按身体区域、压力按取样点。\pz{存疑：上述“$100$ 名受试者、$29$ 取样点、$75$ 百分位、力源自汇集 $188$ 源的独立研究、瞬态倍率 $2$（颅/额/面除外）”均已对照 ISO/TS 15066:2016 Annex A 表 A.2 脚注 a/b/c 核实属实；唯标准自陈系“单一研究”、样本不覆盖全人群，后续工作（如更大样本的疼痛阈值复测）已对部分数值提出修订意见，将来版本可能上调样本量与人群覆盖——具体修订留待现行标准正文为准。}

\begin{insight}
准静态是 15066 的“最坏情形锚点”，整个阈值体系绕它构建。这解释了本章后文一个反直觉的结论：碰撞后\emph{主动卸力}（Comply/Retract）远优于\emph{急停锁位}（Reflex）。因为急停后机器人保持位置，若人恰被夹住，接触力就从“瞬态”退化为“准静态”——撞上了那条最严的红线。安全策略的优劣，本质上是看它把碰撞推向准静态还是瞬态那一端。
\end{insight}

\begin{pitfall}[把“满足 ISO 15066”当成“绝对安全”]
两点务必清醒。其一，阈值来自有限样本的疼痛实验（百人量级），不覆盖全人群、且不针对\emph{尖锐}物体——尖锐工具即使总力在阈值内，局部压力也可能因接触面积极小而严重超标（压力 $=$ 力 $/$ 面积）。其二，15066 是“疼痛/轻伤”阈值，不是“零风险”保证，且其法律地位是技术规范而非正式标准。合规是\emph{必要}条件，不是\emph{充分}的安全证明。
\end{pitfall}

% ════════════════════════════════════════════════════════════════
\section{有效碰撞质量与速度上界}
\label{sec:cs-mass}

\paragraph{动机：碰撞峰值力不由“机器人多重”决定。}
一个常见误解是“机械臂越重撞得越狠”。但自由碰撞的峰值力取决于碰撞瞬间的\emph{动能}与接触\emph{刚度}，而参与碰撞的等效惯量是\emph{构型与方向相关}的——不是机器人的总质量。\cref{eq:os-meff} 已给出沿接触法向 $\mathbf n$ 的\textbf{有效碰撞质量}
\begin{equation}
  m_{\rm eff}=\frac{1}{\mathbf n^\top\boldsymbol\Lambda^{-1}\mathbf n},
  \qquad \boldsymbol\Lambda=(\mathbf J\mathbf M^{-1}\mathbf J^\top)^{-1},
  \label{eq:cs-meff}
\end{equation}
$\boldsymbol\Lambda$ 为操作空间惯量（\cref{ch:arm-dyn}）。同一构型，沿不同法向“感觉重量”不同：沿整条臂的纵向碰撞惯量大、$m_{\rm eff}$ 大、更危险；沿末端横向碰撞惯量小、$m_{\rm eff}$ 小、更安全。Franka Panda 这类协作臂的末端 $m_{\rm eff}$ 约 $2$--$8$\,kg（随构型与方向），远低于其总质量。\pz{存疑：$2$--$8$\,kg 与下文 $\approx4$\,kg 为 Tutorial 给出的某构型示例，具体随平台与位形变化，应由该机器人的 $\boldsymbol\Lambda$ 现场计算。}

\paragraph{人—机串联惯量。}
碰撞是机器人有效质量与人体局部有效质量\emph{串联}的两体碰撞，等效质量为二者的调和（串联）组合：
\begin{equation}
  m_{\rm eq}=\Bigl(\frac{1}{m_{\rm eff}}+\frac{1}{m_{\rm human}}\Bigr)^{-1},
  \label{eq:cs-series}
\end{equation}
$m_{\rm human}$ 是被撞身体部位的有效质量（如手部约 $0.6$\,kg）。因调和组合受较小者支配，$m_{\rm eq}$ 总小于两者中较小的一个——撞在轻部位（手指）时等效惯量很小。

\paragraph{速度上界：从力阈值反解。}
把接触建模为有效质量 $m_{\rm eq}$ 撞上刚度 $k$（皮肤 $+$ 机器人结构串联）的弹簧，能量守恒给出穿透到峰值时的弹性势能等于初始动能，由此峰值力与碰撞速度 $v$ 成正比，反解得\textbf{碰撞速度上界}
\begin{equation}
  v_{\max}=\frac{F_{\max}}{\sqrt{m_{\rm eq}\,k}},
  \label{eq:cs-vmax}
\end{equation}
其中 $F_{\max}$ 取目标身体区域的瞬态力阈值 $F_{\rm tr}$（自由碰撞按瞬态判据）。\cref{eq:cs-vmax} 是 PFL 限速的核心：它把“某身体区域不被伤害”翻译成“TCP 速度不得超过 $v_{\max}$”，且 $v_{\max}$ 随构型（$m_{\rm eff}\to m_{\rm eq}$）、接触刚度 $k$、目标区域（$F_{\max}$）三者变化。

\begin{derivation}[由能量守恒得 $v_{\max}$]
设碰撞瞬间法向速度为 $v$，接触近似为线性弹簧 $f=k\delta$（$\delta$ 为穿透深度）。忽略碰撞短时内的驱动与重力做功，初始动能全部转为弹性势能时达峰值穿透 $\delta_{\max}$：
\[
  \tfrac12 m_{\rm eq}v^2=\tfrac12 k\,\delta_{\max}^2
  \;\Rightarrow\; \delta_{\max}=v\sqrt{m_{\rm eq}/k}.
\]
峰值力 $F_{\rm peak}=k\,\delta_{\max}=v\sqrt{k\,m_{\rm eq}}$，与碰撞速度\emph{成正比}。令 $F_{\rm peak}\le F_{\max}$ 即得 \cref{eq:cs-vmax}。
\end{derivation}

\paragraph{数值示例（撞手背）。}
取末端有效质量 $m_{\rm eff}\approx 4$\,kg、手部 $m_{\rm human}\approx 0.6$\,kg，由 \cref{eq:cs-series} 得 $m_{\rm eq}\approx 0.52$\,kg；接触刚度 $k\approx 1.5\times10^{5}$\,N/m，手背瞬态阈值 $F_{\rm tr}=280$\,N。代入 \cref{eq:cs-vmax}：
\begin{equation}
  v_{\max}=\frac{280}{\sqrt{0.52\times 1.5\times10^{5}}}\approx\frac{280}{279}\approx 1.0\ \text{m/s}.
  \label{eq:cs-vmax-num}
\end{equation}
\pz{存疑：接触刚度 $k\approx150$\,N/mm 与该算例为 Tutorial 给出的代表性数值；皮肤/软组织刚度强非线性且因部位、个体差异巨大，所得 $v_{\max}$ 仅为量级参考，认证须实测或按标准方法评估。}

\paragraph{与阻抗参数的关系——及“不等于自动满足 PFL”。}
若阻抗控制器设刚度 $\mathbf K_d$、误差限幅 $e_{\max}$（如操作空间章的限幅阻抗），则\emph{准静态}最大接触力有简单上界
\begin{equation}
  F_{\rm qs,max}=\lambda_{\max}(\mathbf K_d)\,e_{\max},
  \label{eq:cs-Fqs}
\end{equation}
例如 $\mathbf K_d=500$\,N/m、$e_{\max}=0.05$\,m 给出 $F_{\rm qs,max}=25$\,N，低于多数区域的准静态阈值。但\textbf{这绝不等于自动满足 PFL}，至少有四条限定（\cref{fig:cs-relation}）：
\begin{enumerate}
  \item \textbf{瞬态峰值另算}：自由碰撞的瞬态峰值力 $F_{\rm peak}\approx v\sqrt{k\,m_{\rm eq}}$（见上推导）由碰撞速度与有效质量决定，与 $\mathbf K_d\,e_{\max}$ \emph{无直接关系}——高速碰撞时即使低刚度阻抗也可能超瞬态阈值。\cref{eq:cs-Fqs} 只约束了\emph{稳态被夹}的准静态力。
  \item \textbf{压力受接触面积制约}：阈值有“力”与“压力”两列，小面积硬接触（棱边、尖角）即使\emph{力}达标，\emph{压力}（力/面积）仍可能超标。\cref{eq:cs-Fqs} 完全没碰压力这一维。
  \item \textbf{需逐区域验证}：须对\emph{每个}可能被撞的身体区域（约 $29$ 个）分别核对力与压力两列，而非只看“某个”阈值。
  \item \textbf{需配合限速与碰撞后响应}：完整 PFL 合规是“限速（防瞬态超标）$+$ 检测（本章第~\ref{sec:cs-detection}~节）$+$ 碰撞后主动卸力（第~\ref{sec:cs-postcollision}~节）”的闭环，阻抗参数设计只是其中一环。
\end{enumerate}

\begin{figure}[ht]
\centering
\begin{tikzpicture}[
  >=Stealth, font=\small, node distance=5mm and 9mm,
  src/.style={draw, rounded corners, align=center, inner sep=3pt, fill=main!6, draw=main!60!black},
  chk/.style={draw, rounded corners, align=center, inner sep=3pt, fill=second!6, draw=second!60!black},
  bad/.style={draw, rounded corners, align=center, inner sep=3pt, fill=red!4, draw=red!55!black}]
\node[src] (imp) {阻抗设计\\ $F_{\rm qs,max}=\lambda_{\max}(\mathbf K_d)\,e_{\max}$};
\node[chk, right=of imp] (qs) {准静态力\\ 达标？};
\node[bad, above right=4mm and 10mm of qs] (tr) {瞬态峰值\\ $v\sqrt{k\,m_{\rm eq}}$};
\node[bad, right=of qs] (pr) {压力\\ 力/面积};
\node[bad, below right=4mm and 10mm of qs] (cl) {逐区域 +\\ 限速 + 碰后响应};
\draw[->, thick, gray!60!black] (imp)--(qs);
\draw[->, thick, red!55!black] (qs)--(tr); \draw[->, thick, red!55!black] (qs)--(pr); \draw[->, thick, red!55!black] (qs)--(cl);
\node[font=\scriptsize, red!55!black, align=center] at (8.4,-1.55) {以上任一不满足\\ $\Rightarrow$ 不合规};
\end{tikzpicture}
\caption{“准静态力达标”只是 PFL 合规的一个必要分支。瞬态峰值、压力（接触面积）、逐区域核对与限速/碰后响应均须独立满足——阻抗参数设计不能单独证明 PFL 合规。}
\label{fig:cs-relation}
\end{figure}

\begin{pitfall}[用 $\mathbf K_d\,e_{\max}$ “证明”了安全]
这是力控新手最常见的合规陷阱：算出准静态力 $25$\,N $<$ 阈值，便宣告“满足 ISO 15066”。\cref{fig:cs-relation} 显示这只覆盖了准静态力一支。真正的反例：以 $1$\,m/s 携同一低刚度阻抗自由撞向人手，瞬态峰值力可达数百牛（\cref{eq:cs-vmax-num} 的反向）；或末端带尖锐工具，力 $10$\,N 但接触面积 $0.1$\,cm$^2$、压力 $100$\,N/cm$^2$ 严重超标。准静态达标\emph{必要}但远\emph{不充分}。
\end{pitfall}

% ════════════════════════════════════════════════════════════════
\section{碰撞后策略：Reflex / Retract / Comply}
\label{sec:cs-postcollision}

\paragraph{动机：检测到碰撞之后做什么。}
检测只是第一步。碰撞被判定后，控制器必须立刻\emph{切换模式}收场——而“怎么收场”直接决定了人受到的是瞬态载荷还是准静态载荷。三种典型策略（\cref{fig:cs-strategies}）：
\begin{itemize}
  \item \textbf{Reflex（反射急停）}：立即冻结当前位置（高刚度锁位、零速度指令），把机器人钉在原地。
  \item \textbf{Retract（主动回缩）}：沿碰撞反方向匀速退开一小段，主动\emph{脱离}接触。
  \item \textbf{Comply（柔顺让步）}：切换到零刚度、过阻尼的“零力引导”模式，让接触力把机器人\emph{推开}，人可徒手推动它。
\end{itemize}

\begin{figure}[ht]
\centering
\begin{tikzpicture}[>=Stealth, font=\small, scale=1.0]
  % Reflex
  \node[font=\bfseries] at (-3.6,1.9) {Reflex 急停};
  \fill[second!70!black] (-4.5,0.4) circle (0.16);
  \draw[main!60!black, very thick] (-4.25,1.2) rectangle (-3.85,0.7);
  \draw[->, red!65!black, very thick] (-4.05,0.65) -- (-4.05,0.5);
  \draw[gray!60!black, very thick] (-4.9,-0.05) -- (-3.2,-0.05);
  \node[font=\scriptsize, red!65!black] at (-3.4,0.35) {力滞留};
  \draw[red!65!black] (-3.95,0.95) -- (-3.7,0.95);% locked marker
  \node[font=\scriptsize] at (-3.85,1.45) {锁位};
  % Retract
  \node[font=\bfseries] at (-0.2,1.9) {Retract 回缩};
  \fill[second!70!black] (-1.1,0.4) circle (0.16);
  \draw[main!60!black, very thick] (-0.55,1.2) rectangle (-0.15,0.7);
  \draw[->, third!70!black, very thick] (-0.6,0.95) -- (-1.05,0.95) node[midway,above,font=\scriptsize]{退};
  \draw[gray!60!black, very thick] (-1.5,-0.05) -- (0.4,-0.05);
  \node[font=\scriptsize, third!70!black] at (0.1,0.3) {脱离};
  % Comply
  \node[font=\bfseries] at (3.3,1.9) {Comply 柔顺};
  \fill[second!70!black] (2.4,0.4) circle (0.16);
  \draw[main!60!black, very thick] (2.95,1.2) rectangle (3.35,0.7);
  \draw[<-, main!60!black, very thick] (2.9,0.95) -- (3.4,0.95);
  \draw[->, second!70!black, thick] (2.55,0.55) -- (2.9,0.7) node[right,font=\scriptsize]{人推开};
  \draw[gray!60!black, very thick] (2.0,-0.05) -- (3.9,-0.05);
  \node[font=\scriptsize, main!60!black] at (3.7,0.3) {让步};
\end{tikzpicture}
\caption{三种碰撞后策略。Reflex 锁住位置——若人被夹在机器人与固定面间，接触力\emph{滞留}成准静态载荷（最危险）；Retract 主动退开使接触脱离；Comply 转零力引导，让接触力把机器人推开。}
\label{fig:cs-strategies}
\end{figure}

\paragraph{Reflex 的悖论：残余夹持力。}
直觉上“撞到就刹停锁住”最安全，但在 PFL 框架下 Reflex 恰是三者里\textbf{最危险}的——这是一个值得专门点破的安全悖论。原因正是第~\ref{sec:cs-iso}~节的准静态/瞬态二分：Reflex 急停后机器人\emph{保持位置不动}，若人恰被夹在机器人与固定物体之间，接触力\textbf{无法卸去、持续施加}，于是从“瞬态”载荷退化为“准静态”载荷——而准静态阈值是两套里最严的（约为瞬态的一半，\cref{tab:cs-iso}）。换言之，Reflex 把一次本可短暂卸去的瞬态碰撞，\emph{冻结}成了持续压迫的准静态夹持。\textbf{急停消除了运动，却没有消除力}——这正是它的悖论所在。

\begin{insight}
“停下来”与“卸掉力”是两件不同的事，安全要的是后者。Reflex 实现了前者却放弃了后者：它把动能归零，却让接触势能（弹性变形对应的力）持续保持。Comply/Retract 之所以更优，正因为它们\emph{主动减小接触力}——Comply 撤掉回复刚度让力自然衰减、Retract 直接脱离接触。把这条想透，就理解了为什么协作安全的“反射动作”反直觉地不是“刹车”而是“松手/后退”。
\end{insight}

\paragraph{三策略的合规对比。}
\cref{tab:cs-postcompare} 以“以低速撞上弹性物体”的典型评估口径比较三者。关键看两行：\emph{碰撞后残余力}与\emph{准静态合规}。Reflex 残余力滞留、准静态\textbf{不合规}；Comply/Retract 残余力趋零、准静态合规。

\begin{table}[H]\centering\small
\caption{碰撞后三策略对比（低速撞弹性物体的典型评估口径）}
\label{tab:cs-postcompare}
\begin{tabular}{lccc}
\toprule
指标 & Reflex（急停） & Retract（回缩） & Comply（柔顺） \\
\midrule
峰值碰撞力 & 较高 & 中 & \textbf{最低} \\
碰撞后残余力 & \textbf{高（持续夹持）} & 趋零（已脱离） & 趋零（主动让步） \\
恢复耗时 & 最短（但不安全） & 中 & 中--长 \\
ISO 15066 准静态合规 & \textbf{不合规}（持续力） & 合规 & 合规 \\
实现复杂度 & 低 & 中 & 中 \\
\bottomrule
\end{tabular}
\end{table}

\begin{finenote}
\pz{存疑：\cref{tab:cs-postcompare} 的相对优劣来自 Tutorial 给出的一组 MuJoCo 仿真评估（如以 $0.3$\,m/s 撞 $K_e$ 数千 N/m 的弹性体，Reflex 残余力数十牛、Comply 峰值最低）；具体数值依仿真设置而定，此处只保留\emph{相对}结论（残余力 Reflex 高、Comply/Retract 趋零），绝对数值不作断言。}
\end{finenote}

\paragraph{策略选择：按场景而非按直觉。}
没有一种策略普适，应按碰撞性质选择：
\begin{itemize}
  \item \emph{瞬态/高速碰撞}：先 Reflex 止住动能、再 Retract 脱离——\textbf{先停再退}，避免停后退化成准静态夹持；
  \item \emph{准静态/低速夹持}：直接 Comply，切零力引导让人能把机器人推开——主动卸力是这里唯一安全的选项；
  \item \emph{检测到人进入协作区（尚未接触）}：退回 SSM 的预防性停止（最保守）。
\end{itemize}

\begin{algo}[碰撞后响应状态机（控制级，伪代码）]\label{alg:cs-recovery}
状态机在控制频率下运行，四态依次为正常监测、碰撞确认、策略执行、交还恢复：
\begin{enumerate}
  \item \textbf{NORMAL}：正常控制；持续监测 $\|\mathbf r\|$ 的方向分离分量 $r_{\rm op},r_\perp$（\cref{eq:cs-dirsplit}）。任一越其阈值 $\to$ DETECTED，记下碰撞方向 $\hat{\mathbf n}$ 与当前位姿。
  \item \textbf{DETECTED}：据策略初始化——Reflex 置高刚度锁位、零速度指令；Retract 记回缩起点、置沿 $-\hat{\mathbf n}$ 的退缩速度、撤回复刚度；Comply 置零刚度、过阻尼。转 RECOVERING。
  \item \textbf{RECOVERING}：执行所选策略。Retract 至退缩距离达标即完成；Comply 至外力降到清除阈值（如数牛）即完成；超时（如数秒）强制完成。转 RECOVERED。
  \item \textbf{RECOVERED}：交还正常控制器，由其用平滑过渡（如五次多项式插值刚度，使 $\dot{\mathbf K}_d$ 端点连续，与操作空间章一致）逐步恢复任务参数，避免二次冲击。回 NORMAL。
\end{enumerate}
\end{algo}

\begin{pitfall}[Retract 沿“错误方向”退缩]
Retract 必须沿\emph{碰撞接触法向的反方向} $-\hat{\mathbf n}$ 退缩。若用了过期的运动方向、或在多点接触下取了某条连杆的局部法向，可能退向\emph{另一}障碍、或把被夹的人体往更糟的方向带。回缩方向应取自当前碰撞估计（残差反投影的 wrench 方向），且回缩中持续监测残差——一旦反方向又出现接触力，立即转 Comply。
\end{pitfall}

% ════════════════════════════════════════════════════════════════
\section{能量分级响应与本章小结}
\label{sec:cs-energy}

\paragraph{把检测、限速、碰后策略接成一条链。}
前面各节是离散的工具：检测给“是否碰撞”、限速给“速度上界”、碰后策略给“怎么收场”。运行时还需要一个\emph{连续}的安全裕度度量把它们串起来，并据裕度消耗程度采取\emph{分级}干预——而非“正常/急停”的二值逻辑。一个自然的度量是\textbf{能量}：以端口功率 $P=\mathbf f^\top\mathbf v$（\cref{eq:os-port-power}）累积系统与环境交换的能量，越界则分级响应。这与遥操作中按端口能量流监控无源性的时间域无源性方法（\cref{ch:teleop-tdpa}）共享同一记账思想——彼处用能量赤字注入附加阻尼以保通信通道无源，此处用能量裕度触发分级安全响应。\cref{tab:cs-levels} 给出四级方案。

\begin{table}[H]\centering\small
\caption{能量分级响应（Level 0--3）：按安全裕度消耗程度递增干预}
\label{tab:cs-levels}
\begin{tabular}{llp{6.4cm}}
\toprule
级别 & 触发（裕度消耗） & 干预动作 \\
\midrule
\textbf{Level 0} 正常 & 裕度充足、净能量流低 & 不干预，仅记录日志 \\
\textbf{Level 1} 注意 & 裕度中等或净功率升高 & 低优先级告警；降低最大允许速度/刚度变化率，主动减缓能量消耗 \\
\textbf{Level 2} 警告 & 裕度偏低或净功率持续为正 & 高优先级告警；冻结升刚度与加速，提高阻尼以加速能量耗散 \\
\textbf{Level 3} 紧急 & 裕度耗尽或累积交换能量越上限 & 触发碰撞后策略（首选 Comply 卸力）；必要时安全停机 \\
\bottomrule
\end{tabular}
\end{table}

\begin{insight}
分级响应的价值在于\emph{尽量保住任务连续性的同时不放弃安全}：二值逻辑要么放任、要么急停（而急停又可能引出 Reflex 的准静态悖论）。能量度量天然连续、可微分，使“轻踩刹车”成为可能——这也呼应了第~\ref{sec:cs-postcollision}~节的核心结论：安全的终极动作是\emph{管理能量/力}，而非简单地\emph{停止运动}。Level 3 默认走 Comply 而非 Reflex，正是这一哲学的落点。
\end{insight}

\begin{pitfall}[能量监测的采样率不足]
能量裕度由功率 $P=\mathbf f^\top\mathbf v$ 在每个控制周期累积。若把监测放在低频线程（如 $10$\,Hz），两次采样之间的高频力脉冲会被\emph{漏积}——一次短促碰撞的能量未被记入，分级响应失灵。能量监测必须与力控同频（典型 $1$\,kHz）、在同一线程内计算，方与检测/控制的时间分辨率一致。
\end{pitfall}

\paragraph{本章小结。}
本章把机械臂安全从“能写力控律”推进到“能让与人共处的机械臂在碰撞时不伤人、且事后安全收场”，核心是把 ISO/TS 15066 这条法规红线翻译成一串控制约束：
\begin{enumerate}
  \item \textbf{两种模式分工}：SSM 靠感知限速预防接触、力控不直接守红线；PFL 允许接触、由力控把接触力压在阈值内——本章聚焦 PFL（\cref{sec:cs-motivation}）。
  \item \textbf{检测工程化}：以 \cref{thm:os-observer} 的残差 $\mathbf r$ 为输入，阈值取略高于噪声底，并用方向分离 $r_{\rm op}/r_\perp$（\cref{eq:cs-dirsplit}）在“放过工艺力”与“抓住侧向异常”间取灵敏度（\cref{sec:cs-detection}）。
  \item \textbf{合规阈值}：15066 按约 $29$ 区域、区分准静态/瞬态给出力/压力阈值（\cref{tab:cs-iso}），瞬态$\approx 2\times$准静态源于组织对短时载荷耐受更高的生物力学（\cref{sec:cs-iso}）。
  \item \textbf{速度上界}：有效质量 $m_{\rm eff}=1/(\mathbf n^\top\boldsymbol\Lambda^{-1}\mathbf n)$ 构型相关，由能量守恒得 $v_{\max}=F_{\max}/\sqrt{m_{\rm eq}k}$；准静态力 $\lambda_{\max}(\mathbf K_d)e_{\max}$ 达标\emph{必要不充分}，瞬态峰值、压力、逐区域、限速/碰后响应均须独立满足（\cref{sec:cs-mass}）。
  \item \textbf{碰撞后策略}：Reflex 急停会把瞬态碰撞冻结成准静态夹持（残余夹持力悖论），Comply/Retract 主动卸力故准静态合规、应为首选（\cref{sec:cs-postcollision}）。
  \item \textbf{运行时链}：以能量为度量做 Level\,0--3 分级响应，把检测、限速、碰后策略接成一条连续的安全链，紧急级默认 Comply 卸力（\cref{sec:cs-energy}）。
\end{enumerate}

\begin{practice}
\begin{enumerate}
  \item \textbf{速度上界表}：给定某协作臂末端有效质量 $4$\,kg、接触刚度 $1.5\times10^{5}$\,N/m，对 \cref{tab:cs-iso} 中胸骨、手背、大腿三区域，用 \cref{eq:cs-series,eq:cs-vmax}（人体局部质量自设合理值、$F_{\max}$ 取各区瞬态力）算各自的 $v_{\max}$；再说明：为何\emph{前额}一行不能照此算出一个有限 $v_{\max}$（提示：标准对颅/额/面不给瞬态值、接触不被允许），以及这如何对应“头面部限速最严”的极端形式。
  \item \textbf{合规反例}：阻抗控制器取 $\mathbf K_d=1000$\,N/m、$e_{\max}=0.05$\,m。其准静态最大力是多少？满足 \cref{tab:cs-iso} 所有区域的准静态阈值吗？再构造一个该参数下\emph{瞬态}超标的碰撞场景（给出速度与有效质量），说明 \cref{eq:cs-Fqs} 为何不充分。
  \item \textbf{方向分离设计}：在销孔装配插入阶段，沿 $z$ 的工艺力可达 $30$\,N、侧向卡死信号约 $5$\,N。按 \cref{eq:cs-dirsplit} 给 $r_{\rm op}$ 与 $r_\perp$ 各设一个阈值，使工艺力不误触发、侧向 $5$\,N 能被抓住；再说明为什么单一标量阈值无法两全。
  \item \textbf{策略论证}：用准静态/瞬态二分（\cref{sec:cs-iso}）论证“人被夹在机器人与桌面之间”时为何 Comply 优于 Reflex；若改为“机器人在自由空间高速擦碰到人后人已离开”，Reflex$+$Retract 是否更合适？为什么？
\end{enumerate}
\end{practice}

\begin{table}[H]\centering\small
\caption{本章符号表}
\label{tab:acs-symbols}
\begin{tabular}{ll}
\toprule
符号 & 含义 \\
\midrule
$\mathbf r$ & 动量观测器残差，估计外力矩 $\boldsymbol\tau_{\rm ext}$（\cref{thm:os-observer}） \\
$r_{\rm th}$ & 碰撞检测阈值（残差范数） \\
$r_{\rm op},\,r_\perp$ & 残差/wrench 沿操作方向与正交补的分量（方向分离，\cref{eq:cs-dirsplit}） \\
$\hat{\mathbf n}_{\rm op},\,\hat{\mathbf n}$ & 单位操作方向 / 碰撞接触法向 \\
$F_{\rm qs},\,F_{\rm tr}$ & ISO/TS 15066 准静态 / 瞬态力阈值 \\
$p_{\rm qs},\,p_{\rm tr}$ & 准静态 / 瞬态压力阈值 \\
$m_{\rm eff}$ & 构型相关有效碰撞质量 $1/(\mathbf n^\top\boldsymbol\Lambda^{-1}\mathbf n)$（\cref{eq:cs-meff}） \\
$m_{\rm human},\,m_{\rm eq}$ & 被撞身体部位有效质量 / 人—机串联等效质量（\cref{eq:cs-series}） \\
$\boldsymbol\Lambda$ & 操作空间惯量 $(\mathbf J\mathbf M^{-1}\mathbf J^\top)^{-1}$（\cref{ch:arm-dyn}） \\
$k$ & 接触刚度（皮肤 $+$ 机器人结构串联） \\
$v_{\max}$ & PFL 碰撞速度上界 $F_{\max}/\sqrt{m_{\rm eq}k}$（\cref{eq:cs-vmax}） \\
$\mathbf K_d,\,\mathbf D_d,\,e_{\max}$ & 阻抗刚度 / 阻尼 / 误差限幅 \\
$P=\mathbf f^\top\mathbf v$ & 接触端口功率（\cref{eq:os-port-power}），分级响应的能量度量 \\
\bottomrule
\end{tabular}
\end{table}
